%!TEX root = ../../main.tex
%% ===============================================================
%% 4.3 공간 지름 분할
%% 파일: chapters/04_localization/03_spatial_split.tex
%% 상위: chapters/04_localization/chapter.tex
%% 정의 라벨: sec:loc-spatial
%% 참조 라벨: def:engagement, sec:disc-audit
%% 포함 객체: -
%% 인용: defays1977efficient, sorensen1948method
%% 상태: 초안 (docs/OUTLINE.md 참고)
%% ===============================================================

\section{공간 지름 분할}
\label{sec:loc-spatial}

시간 군집은 서로 다른 장소에서 동시에 벌어진 교전(예: 탑 라인 소규모 교전과 봇 라인 다이브)을 하나로 묶을 수 있다. 이를 분리하기 위해 군집의 \emph{공간 지름}
\begin{equation}
  \mathrm{diam}(C) = \max_{u, u' \in C} \big\|\mathrm{pos}(u) - \mathrm{pos}(u')\big\|_2
\end{equation}
이 $D_{\max} = 4{,}000$ 단위를 넘으면 지름이 $D_{\max}$ 이하인 부분 군집들로 나눈다. 이는 완전 연결(complete-linkage) 기준의 응집적 군집화 \cite{sorensen1948method,defays1977efficient}에 해당한다. 단일 연결을 쓰면 킬이 사슬처럼 이어질 때 총 지름이 $D_{\max}$를 초과할 수 있으므로, 완전 연결이 정의 \ref{def:engagement}의 ``지름 상한''을 정확히 구현한다.\footnote{초기 구현은 단일 연결이었으며, 사전 감사에서 완전 연결로 수정되었다(제 \ref{sec:disc-audit} 절).} 군집의 모든 쌍 거리가 $D_{\max}$ 이하이면 분할 없이 통과하는 빠른 경로를 둔다.
